\chapter{Retas, círculos e ângulos}\label{ch:g6:lines}

A geometria começa com régua, esquadro e compasso. Este capítulo fixa o
vocabulário --- pontos, \omterm{def:g6:lines:objects}{segmentos}, semirretas, retas, círculos ---,
apresenta as duas posições especiais de retas (\omterm{def:g6:lines:perp}{paralelas} e
\omterm{def:g6:lines:perp}{perpendiculares}) e ensina a medir e a traçar ângulos.

\section{Pontos, segmentos, retas}

\begin{definition}[Objetos básicos]\label{def:g6:lines:objects}
Por dois pontos distintos $A$ e $B$ passam:
\begin{itemize}
  \item o \emph{segmento}\index{segmento} $[AB]$: a parte da reta entre
        $A$ e $B$ (tem um comprimento, escrito $AB$);
  \item a \emph{semirreta} $[AB)$: começa em $A$, passa por $B$ e
        continua para sempre;
  \item a \emph{reta}\index{reta} $(AB)$: prolonga-se para sempre nos
        dois sentidos. Dois pontos determinam exatamente uma reta.
\end{itemize}
Pontos que estão numa mesma reta são chamados de \emph{colineares}.
\end{definition}

\begin{omfigure}
\begin{tikzpicture}[scale=1.0]
  \fill (0,1.8) circle (1.5pt) node[above, font=\small] {$A$};
  \fill (2.5,1.8) circle (1.5pt) node[above, font=\small] {$B$};
  \draw[thick, omDef] (0,1.8) -- (2.5,1.8);
  \node[right, font=\small] at (2.9,1.8) {segmento $[AB]$};
  \fill (0,0.9) circle (1.5pt) node[above, font=\small] {$A$};
  \fill (2.5,0.9) circle (1.5pt) node[above, font=\small] {$B$};
  \draw[thick, omThm] (0,0.9) -- (4.1,0.9);
  \node[right, font=\small] at (4.3,0.9) {semirreta $[AB)$};
  \fill (0,0) circle (1.5pt) node[above, font=\small] {$A$};
  \fill (2.5,0) circle (1.5pt) node[above, font=\small] {$B$};
  \draw[thick, omMeth] (-1.4,0) -- (4.1,0);
  \node[right, font=\small] at (4.3,0) {reta $(AB)$};
\end{tikzpicture}

{\small Os mesmos dois pontos, três objetos diferentes. Os parênteses e
colchetes dizem o que para e o que continua: $[$ para, $($ continua.}
\end{omfigure}

\section{Retas paralelas e perpendiculares}

\begin{definition}[Perpendiculares, paralelas]\label{def:g6:lines:perp}
Duas retas são \emph{perpendiculares}\index{retas perpendiculares}
quando se cruzam formando um \omterm{def:g3:shapes:rightangle}{ângulo reto}; escrevemos
$d_1 \perp d_2$ e marcamos o \omterm{def:g3:shapes:rightangle}{ângulo reto} com um quadradinho. Duas retas
são \emph{paralelas}\index{retas paralelas} quando nunca se encontram,
por mais que sejam prolongadas; escrevemos $d_1 \parallel d_2$.
\end{definition}

\begin{omfigure}
\begin{tikzpicture}[scale=1.0]
  \draw[thick] (-1.8,0) -- (1.8,0) node[right, font=\small] {$d_1$};
  \draw[thick, omDef] (0.6,-1.2) -- (0.6,1.4)
    node[above, font=\small] {$d_2$};
  \draw[thin] (0.6,0) ++(-0.25,0) -- ++(0,0.25) -- ++(0.25,0);
  \node[below, font=\small] at (0,-1.5) {$d_1 \perp d_2$};
  \begin{scope}[xshift=6.5cm]
  \draw[thick] (-1.8,-0.5) -- (1.8,0.3) node[right, font=\small] {$d_3$};
  \draw[thick, omThm] (-1.8,-1.3) -- (1.8,-0.5)
    node[right, font=\small] {$d_4$};
  \node[below, font=\small] at (0,-1.5) {$d_3 \parallel d_4$};
  \end{scope}
\end{tikzpicture}

{\small Um cruzamento em \omterm{def:g3:shapes:rightangle}{ângulo reto} (marcado com o quadradinho) e duas
\omterm{def:g6:lines:perp}{paralelas}: mesma direção, nenhum ponto de cruzamento.}
\end{omfigure}

\begin{proposition}[Dois fatos úteis]\label{prop:g6:lines:facts}
\begin{enumerate}
  \item Se duas retas são ambas \omterm{def:g6:lines:perp}{perpendiculares} a uma terceira reta,
        elas são \omterm{def:g6:lines:perp}{paralelas} entre si.
  \item Se duas retas são \omterm{def:g6:lines:perp}{paralelas}, toda reta perpendicular a uma é
        perpendicular à outra.
\end{enumerate}
\end{proposition}
\admitted

\begin{method}[Traçar com o esquadro]\label{met:g6:lines:setsquare}
Para traçar a perpendicular a uma reta $d$ passando por um ponto $P$:
\begin{enumerate}
  \item apoie uma das bordas do \omterm{def:g3:shapes:rightangle}{ângulo reto} do esquadro sobre $d$;
  \item deslize o esquadro ao longo de $d$ até a outra borda alcançar
        $P$;
  \item trace a reta ao longo dessa borda e marque o \omterm{def:g3:shapes:rightangle}{ângulo reto}.
\end{enumerate}
Para uma paralela por $P$: trace uma perpendicular a $d$ e depois a
perpendicular a \emph{essa} reta passando por $P$
(\cref{prop:g6:lines:facts}).
\end{method}

\section{Círculos}

\begin{definition}[Círculo]\label{def:g6:lines:circle}
O \emph{círculo}\index{círculo} de centro $O$ e \omterm{def:g3:shapes:shapes}{raio} $r$ é o conjunto
de todos os pontos que estão à distância exatamente $r$ de $O$. Um
\omterm{def:g6:lines:objects}{segmento} do centro até o círculo é um \emph{\omterm{def:g3:shapes:shapes}{raio}}; um \omterm{def:g6:lines:objects}{segmento} que liga
dois pontos do círculo passando pelo centro é um \emph{\omterm{def:g4:geometry:circle}{diâmetro}} --- o
comprimento dele é $2r$; um \omterm{def:g6:lines:objects}{segmento} que liga dois pontos do círculo é
uma \emph{corda}.
\end{definition}

\begin{omfigure}
\begin{tikzpicture}[scale=1.15]
  \draw[thick] (0,0) circle (1.5);
  \fill (0,0) circle (1.5pt) node[below left, font=\small] {$O$};
  \fill (1.5,0) circle (1.5pt) node[right, font=\small] {$M$};
  \draw[omDef, thick] (0,0) -- (1.5,0)
    node[midway, above, font=\small] {raio};
  \fill (-1.06,1.06) circle (1.5pt) node[above left, font=\small] {$A$};
  \fill (1.06,-1.06) circle (1.5pt) node[below right, font=\small] {$B$};
  \draw[omThm, thick] (-1.06,1.06) -- (1.06,-1.06)
    node[pos=0.25, right=2pt, font=\small] {diâmetro};
  \fill (0.52,1.41) circle (1.5pt) node[above, font=\small] {$C$};
  \fill (1.41,0.52) circle (1.5pt) node[right, font=\small] {$D$};
  \draw[omMeth, thick] (0.52,1.41) -- (1.41,0.52)
    node[midway, above right=-2pt, font=\small] {corda};
\end{tikzpicture}

{\small Um círculo com um \omterm{def:g3:shapes:shapes}{raio} $[OM]$, um \omterm{def:g4:geometry:circle}{diâmetro} $[AB]$ (uma corda que
passa pelo centro, com o dobro do comprimento do \omterm{def:g3:shapes:shapes}{raio}) e uma corda
$[CD]$.}
\end{omfigure}

\begin{example}\label{ex:g6:lines:circle}
``Trace o círculo de centro $O$ que passa por $A$'': abra o compasso de
$O$ até $A$ --- o \omterm{def:g3:shapes:shapes}{raio} é a distância $OA$ --- e gire. Todo ponto desse
círculo está à mesma distância de $O$ que $A$.
\end{example}

\section{Ângulos}

\begin{definition}[Ângulo]\label{def:g6:lines:angle}
Duas semirretas $[AB)$ e $[AC)$ com o mesmo ponto de partida $A$ formam
o \emph{ângulo}\index{ângulo} $\widehat{BAC}$; o ponto $A$ é o
\emph{\omterm{def:g5:solids:def}{vértice}} dele (sempre a letra do meio!). Os ângulos são medidos
em \emph{graus} ($^\circ$), de $0^\circ$ a $360^\circ$ para uma volta
completa. Um ângulo é:
\begin{itemize}
  \item \emph{reto} se mede $90^\circ$;
  \item \emph{agudo} se mede menos que $90^\circ$;
  \item \emph{obtuso} se mede entre $90^\circ$ e $180^\circ$;
  \item \emph{raso} se mede $180^\circ$ (as duas semirretas formam uma
        reta).
\end{itemize}
\end{definition}

\begin{omfigure}
\begin{tikzpicture}[scale=0.9]
  \draw[thick] (0,0) -- (1.8,0);
  \draw[thick] (0,0) -- (40:1.8);
  \draw[omDef, ->] (0.7,0) arc (0:40:0.7);
  \node[omDef, font=\small] at (20:1.0) {$40^\circ$};
  \node[below, font=\small] at (0.9,-0.3) {agudo};
  \begin{scope}[xshift=3.6cm]
  \draw[thick] (0,0) -- (1.8,0);
  \draw[thick] (0,0) -- (0,1.8);
  \draw[thin] (0.3,0) -- (0.3,0.3) -- (0,0.3);
  \node[below, font=\small] at (0.9,-0.3) {reto};
  \end{scope}
  \begin{scope}[xshift=7.2cm]
  \draw[thick] (0,0) -- (1.8,0);
  \draw[thick] (0,0) -- (135:1.8);
  \draw[omThm, ->] (0.6,0) arc (0:135:0.6);
  \node[omThm, font=\small] at (65:1.0) {$135^\circ$};
  \node[below, font=\small] at (0.9,-0.3) {obtuso};
  \end{scope}
  \begin{scope}[xshift=11.6cm]
  \draw[thick] (-1.5,0) -- (1.5,0);
  \fill (0,0) circle (1.2pt);
  \draw[omMeth, ->] (0.5,0) arc (0:180:0.5);
  \node[below, font=\small] at (0,-0.3) {raso};
  \end{scope}
\end{tikzpicture}

{\small As quatro famílias de ângulos. O \omterm{def:g3:shapes:rightangle}{ângulo reto} ($90^\circ$) é a
referência: agudo quer dizer menor; obtuso, maior.}
\end{omfigure}

\begin{method}[Medir um ângulo com o transferidor]\label{met:g6:lines:protractor}
\begin{enumerate}
  \item Ponha o centro do transferidor exatamente sobre o \omterm{def:g5:solids:def}{vértice} do
        ângulo;
  \item alinhe a linha de $0^\circ$ dele com um dos lados do ângulo;
  \item leia a graduação atravessada pelo outro lado --- usando a
        escala que começa em $0$ no lado alinhado;
  \item confira com o olho: um ângulo agudo tem de dar menos que
        $90^\circ$, e um obtuso, mais.
\end{enumerate}
\end{method}

\begin{example}\label{ex:g6:lines:estimate}
Antes de medir, estime! Um ângulo um pouco mais aberto que o canto de
uma folha de papel tem pouco mais de $90^\circ$; a \omterm{def:g3:division:half}{metade} de um \omterm{def:g3:shapes:rightangle}{ângulo
reto} é $45^\circ$; um terço de um \omterm{def:g3:shapes:rightangle}{ângulo reto} é $30^\circ$. Se o seu
transferidor diz $150^\circ$ para um ângulo que parece agudo, você leu a
escala errada: a medida correta é $180^\circ - 150^\circ = 30^\circ$.
\end{example}

\section{Exercícios}

\begin{exercise}[$\star$]\label{exo:g6:lines:1}
Marque três pontos $A$, $B$, $C$ não colineares. Trace com cores
diferentes: o \omterm{def:g6:lines:objects}{segmento} $[AB]$, a semirreta $[CA)$ e a reta $(BC)$.
\end{exercise}

\begin{exercise}[$\star$]\label{exo:g6:lines:2}
Verdadeiro ou falso? ``$[AB]$ e $[BA]$ são o mesmo \omterm{def:g6:lines:objects}{segmento}.''
``$[AB)$ e $[BA)$ são a mesma semirreta.'' ``$(AB)$ e $(BA)$ são a
mesma reta.'' Explique cada resposta.
\end{exercise}

\begin{exercise}[$\star$]\label{exo:g6:lines:3}
Trace uma reta $d$ e um ponto $P$ fora de $d$. Construa com o esquadro:
a perpendicular a $d$ por $P$ e depois a paralela a $d$ por $P$.
Descreva os seus passos.
\end{exercise}

\begin{exercise}[$\star$]\label{exo:g6:lines:4}
As retas $a$ e $b$ são ambas \omterm{def:g6:lines:perp}{perpendiculares} a uma reta $c$, e uma
quarta reta $e$ é perpendicular a $a$. O que você pode dizer sobre $a$
e $b$? E sobre $e$ e $c$? Justifique com a
\cref{prop:g6:lines:facts}.
\end{exercise}

\begin{exercise}[$\star$]\label{exo:g6:lines:5}
Trace um círculo de centro $O$ e \omterm{def:g3:shapes:shapes}{raio} $3$ cm. Marque um ponto $M$ sobre
o círculo, um ponto $N$ dentro e um ponto $P$ fora. O que você pode
dizer sobre as distâncias $OM$, $ON$, $OP$ em comparação com $3$ cm?
\end{exercise}

\begin{exercise}[$\star$]\label{exo:g6:lines:6}
Um círculo tem \omterm{def:g4:geometry:circle}{diâmetro} $9$ cm. Qual é o \omterm{def:g3:shapes:shapes}{raio} dele? Outro tem \omterm{def:g3:shapes:shapes}{raio}
$2.6$ cm: qual é o \omterm{def:g4:geometry:circle}{diâmetro}?
\end{exercise}

\begin{exercise}[$\star$]\label{exo:g6:lines:7}
Nomeie o ângulo marcado com três letras e depois classifique-o (agudo,
reto, obtuso, raso): um ângulo de $72^\circ$ com \omterm{def:g5:solids:def}{vértice} $R$ entre
semirretas para $S$ e $T$; um ângulo de $148^\circ$ com \omterm{def:g5:solids:def}{vértice} $B$
entre semirretas para $A$ e $C$; um ângulo de $90^\circ$ com \omterm{def:g5:solids:def}{vértice}
$O$ entre semirretas para $M$ e $N$.
\end{exercise}

\begin{exercise}[$\star$]\label{exo:g6:lines:8}
Estime e depois meça com o transferidor os três ângulos de um triângulo
traçado por você. Some as três medidas: o que você encontra? (Guarde a
sua resposta para o \cref{ch:g7:triangles}.)
\end{exercise}

\begin{exercise}[$\star$]\label{exo:g6:lines:9}
Trace um ângulo $\widehat{xOy}$ de $65^\circ$ com o transferidor e
depois um ângulo de $130^\circ$. Como você conseguiria o segundo a
partir do primeiro sem o transferidor?
\end{exercise}

\begin{exercise}[$\star\star$]\label{exo:g6:lines:10}
Dois vilarejos $A$ e $B$ estão marcados num mapa. Onde ficam os pontos
que estão a $3$ cm de $A$ \emph{e} a $2$ cm de $B$? Faça um desenho
mostrando quantos pontos desses pode haver ($0$, $1$ ou $2$, conforme a
distância $AB$).
\end{exercise}

\begin{exercise}[$\star\star$]\label{exo:g6:lines:11}
Um relógio marca 3 horas: qual é o ângulo entre os dois ponteiros? A
mesma pergunta às 5 horas e às 6 horas. (Uma volta completa é
$360^\circ$ para $12$ horas.)
\end{exercise}

\begin{exercise}[$\star\star\star$]\label{exo:g6:lines:12}
Trace um \omterm{def:g6:lines:objects}{segmento} $[AB]$ de $6$ cm. Construa o ponto $C$ tal que
$AC = BC = 6$ cm, usando apenas o compasso, e meça o ângulo
$\widehat{CAB}$. Que triângulo você construiu, e o que você conjectura
sobre os ângulos dele?
\end{exercise}

\section{Problema: a geometria do mostrador do relógio}

\begin{problem}[{Problema de fim de semana --- ângulos como frações de
volta: lê-los no relógio, caçar os momentos em que os ponteiros se
encontram e fatiar um dia numa pizza}]
\label{pb:g6:lines:1}
O relógio é um transferidor que diz as horas: o mostrador dele é uma
volta completa de $360^\circ$, cortada pelas doze marcas das horas em
doze ângulos iguais. O \cref{exo:g6:lines:11} mediu os ponteiros às
3 horas e às 6 horas; este problema constrói a teoria completa ---
incluindo os horários em que \emph{nenhum} ponteiro aponta para uma
marca ---, responde a uma pergunta que poucos adultos acertam
(``quantas vezes os dois ponteiros ficam exatamente um sobre o
outro?'') e termina fatiando um dia inteiro num gráfico de pizza.

\medskip
\noindent\textbf{Parte I --- Frações de volta.}
\begin{enumerate}
  \item Uma volta completa mede $360^\circ$. Quantos graus são meia
        volta, um quarto de volta, um doze avos de volta?
  \item As doze marcas das horas cortam o mostrador em doze ângulos
        iguais no centro. Quantos graus entre duas marcas vizinhas?
        Recupere a resposta do \cref{exo:g6:lines:11} para as 3 horas
        contando marcas.
  \item Dê o ângulo entre os ponteiros à 1 hora, às 4 horas e às
        7 horas. (Às 7 horas os ponteiros separam o mostrador em dois
        ângulos; ``o ângulo entre os ponteiros'' sempre significa o
        \emph{menor} deles.)
  \item O ponteiro dos minutos dá uma volta completa em $60$ minutos.
        Quantos graus ele varre por minuto?
  \item O ponteiro das horas vai de uma marca à seguinte ---
        $30^\circ$ --- em $60$ minutos. Quantos graus \emph{ele} varre
        por minuto? (Um \omterm{def:g6:decimals:places}{número decimal},
        \cref{ch:g6:decimals}.)
\end{enumerate}

\medskip
\noindent\textbf{Parte II --- Horários em que nada aponta para uma
marca.}
\begin{enumerate}[resume]
  \item Às 3:30, o ponteiro dos minutos aponta para o $6$. Onde está
        exatamente o ponteiro das horas? Calcule o ângulo de cada
        ponteiro a partir do $12$ (medindo no sentido horário) e deduza
        o ângulo entre os ponteiros.
  \item Às 6:30, muita gente acha que os ponteiros estão um sobre o
        outro. Chute primeiro e depois calcule o ângulo como na questão
        6. Quem estava certo?
  \item Calcule o ângulo entre os ponteiros às 9:15.
  \item Explique por que, em algum instante entre 1:00 e 1:10, os dois
        ponteiros têm de estar exatamente um sobre o outro: onde está o
        ponteiro dos minutos em relação ao das horas à 1:00, e onde
        está à 1:10? Que ponteiro corre mais depressa, e por que isso
        resolve a questão?
  \item Em doze horas, quantas vezes os dois ponteiros ficam exatamente
        um sobre o outro? (Um encontro acontece em cada trecho entre
        horas sucessivas --- com uma exceção: o que acontece entre 11 e
        12? Liste os horários aproximados dos encontros e conte.) E num
        dia inteiro?
\end{enumerate}

\medskip
\noindent\textbf{Parte III --- Ângulos em torno de um ponto: gráficos
de pizza.}
\begin{enumerate}[resume]
  \item Os doze setores do mostrador juntos preenchem o mostrador: os
        ângulos deles somam $360^\circ$. Explique por que isso vale
        para \emph{qualquer} coleção de ângulos que compartilhem um
        \omterm{def:g5:solids:def}{vértice} e preencham uma volta completa em torno dele, sem
        sobreposição e sem falha.
  \item Zoe registra o dia dela, de $24$ horas: sono $9$ h, escola
        $6$ h, brincadeira $3$ h, refeições $2$ h, todo o \omterm{def:g3:division:remainder}{resto} $4$ h.
        Ela quer um \emph{gráfico de pizza}: um disco em que cada
        atividade ganha um setor, com ângulos proporcionais aos tempos.
        Que \omterm{def:g6:fractions:def}{fração} do dia é cada atividade, e quantos graus ganha o
        setor dela? Confira que os cinco ângulos somam $360^\circ$.
  \item Descreva, passo a passo, como traçar o gráfico de pizza de Zoe
        com compasso e transferidor
        (\cref{met:g6:lines:protractor}): para onde vai o primeiro
        \omterm{def:g3:shapes:shapes}{raio} e como cada setor novo começa onde o anterior termina.
  \item Em outro gráfico de pizza de um dia, um setor mede exatamente
        $90^\circ$. Que \omterm{def:g6:fractions:def}{fração} do disco é essa, e quantas horas ela
        representa?
  \item O grande final, de volta ao relógio: em $24$ horas, quantas
        voltas completas dá o ponteiro das horas? E o dos minutos? E o
        dos segundos? (Uma dessas respostas passa de mil.)
\end{enumerate}
\end{problem}
